% Nouveau EGA V (cote 25): archive cote 25, site page 5 (PDF page 6).
% typescript. Draft compared: none.
% AI-transcribed by Claude Opus 5.5 (claude-opus-5-5, adaptive thinking, effort high), Message Batches API, batch ega25b (msgbatch_01A7n5YuXubUPNK1zsepi8RG), 2026-11-13; not yet checked by eye.
\scanpage{25/5}
\mspage{2}

ou encore $\underline{L}^{-1}\subset\underline{E}$) et $\emptyset$ s'obtient en prenant\var{\inserted{pour $\emptyset$}}{pour $\emptyset$}{insertion}{ega25_25_005-v0001} une section de $\underline{L}^{-1}$ qui
ne s'annule en aucun point \var{\xout{\illegible{1}}}{}{typewriter strike-out}{ega25_25_005-v0002}, i.e. une trivialisation de $\underline{L}^{-1}$
(qui existe en tous cas localement). Notons d'ailleurs que $H_{\hand{\xi}}$ n'est
autre que $P(\underline{E}/\underline{L}^{-1})$ (isomorphisme canonique), ce qui est une troisième
façon de décrire $H_{\hand{\xi}}$ (NB $P(\underline{E}/\underline{L}^{-1})$ est bien plongé canoniquement dans
$P = P(\underline{E})$ ), qui a l'avantage de montrer de plus que $H$ est un fibré projec-%
tif sur $S$, et à fortiori est lisse sur $S$ (encore aurait-il fallu dire \var{\xout{\unsure{qu}}}{}{typewriter strike-out}{ega25_25_005-v0003}
\var{\margin{$\|$}}{\placedmargin{$\|$}}{marginal note placed}{ega25_25_005-v0004} au par. 17 qu'un fibré projectif est lisse ...). Il serait sans doute
\var{\margin{$\|$}}{\placedmargin{$\|$}}{marginal note placed}{ega25_25_005-v0005} préférable de commencer par celle-là.
\uline{Remarques}. La formation de $H_{\hand{\xi}}$ en termes de $\hand{\xi}$ est compatible avec
le changement de base, comme on constate aussitôt., en d'autres termes,
on trouve un homomorphisme de foncteurs $(\mathrm{Sch})^{\circ}_{/S}\to(\mathrm{Ens})$
\[ \check{P} \longrightarrow \underline{\mathrm{Div}}\,(P/S) \qquad , \]
où le deuxième membre désigne le foncteur des ``diviseurs relatifs'' de $P/S$,
dont la valeur en $S'$ ($S$-préschéma arbitraire) est l'ensemble des
sous-schémas fermés de $P_{S'}$ qui sont intersection complète transversale
de codimension 1 relativement à $S'$ (cf par.19). Il est facile de montrer
que cet\var{\xout{\illegible{1}}}{}{typewriter strike-out}{ega25_25_005-v0006} homomorphisme de foncteurs est un monomorphisme, en d'autres
termes que $\hand{\xi}$ est déterminé par $H_{\hand{\xi}}$ . (Ce dernier fait justifie la termino-%
logie ``schéma des hyperplans'' utilisée plus \var{\xout{\unsure{a}}}{}{typewriter strike-out}{ega25_25_005-v0007}haut). Nous verrons \var{\xout{\illegible{1}}}{}{typewriter strike-out}{ega25_25_005-v0008} \var{\margin{)}}{\placedmargin{)}}{marginal note placed}{ega25_25_005-v0009}
\var{\xout{\unsure{Chap.V}}}{}{typewriter strike-out}{ega25_25_005-v0010} que le foncteur $\underline{\mathrm{Div}}$ $(P/S)$ est représentable par \var{\xout{\unsure{un}}}{}{typewriter strike-out}{ega25_25_005-v0011}\var{\inserted{le pré}}{le pré}{insertion}{ega25_25_005-v0012}schéma somme
des $P(\mathrm{Symm}^{k}(\check{E}))$, de sorte que $\overset{\hand{\vee}}{P}$ s'identifie à un sous-préschéma ouvert
et fermé de $\underline{\mathrm{Div}}$ $(P/S)$ ... (NB \var{\over{A}{À}}{À}{writing over}{ega25_25_005-v0013} vrai dire, la d\var{\over{e}{é}}{é}{writing over}{ega25_25_005-v0014}termination des diviseurs
\var{\margin{$|$}}{\placedmargin{$|$}}{marginal note placed}{ega25_25_005-v0015} relatifs\var{\inserted{de $P/S$}}{de $P/S$}{insertion}{ega25_25_005-v0016} pourrait se faire avec les moyens du bord dès maintenant, en
\var{\margin{$|$}}{\placedmargin{$|$}}{marginal note placed}{ega25_25_005-v0017} utilisant les résultats du Chap III, et pourrait être ajoutée \var{\xout{\illegible{1}}}{}{typewriter strike-out}{ega25_25_005-v0018} à titre
\var{\margin{$|$}}{\placedmargin{$|$}}{marginal note placed}{ega25_25_005-v0019} d'exemple au par.19 ...).

\mspage{5}
